\frfilename{mt211.tex}
\versiondate{20.11.03}
     
\def\chaptername{Taksonomi ruang ukur}
\def\sectionname{Definisi}
     
\newsection{211}
     
Saya mulai dengan suatu daftar definisi, yang bersesuaian dengan
konsep-konsep yang menurut pengalaman saya berguna untuk membedakan
berbagai jenis ruang ukur.   Tak terelakkan, arti penting banyak gagasan
ini mungkin masih belum jelas sampai Anda menjumpai beberapa kendala
yang muncul kemudian.   Meskipun demikian, saya berharap Anda akan
dapat menangani definisi-definisi ini secara formal dan abstrak, serta
mengikuti argumen elementer yang digunakan untuk menetapkan hubungan
di antara mereka (211L).
     
Dalam 216C-216E %216C 216D 216E
di bawah ini saya akan memberikan tiga contoh substansial untuk
memperlihatkan betapa beragamnya objek yang tercakup oleh definisi
`ruang ukur'.   Oleh karena itu, dalam bagian ini saya cukup
memberikan uraian yang sangat singkat mengenai sejumlah kasus yang
memadai untuk setidaknya menunjukkan bahwa setiap definisi di sini
membedakan ruang-ruang yang berlainan (211M-211R).
     
\leader{211A}{Definisi} Misalkan $(X,\Sigma,\mu)$ suatu ruang ukur.
Maka $\mu$, atau $(X,\Sigma,\mu)$, disebut ({\bf \Caratheodory}) {\bf
lengkap} jika, setiap kali $A\subseteq E\in\Sigma$ dan $\mu E=0$,
berlaku $A\in\Sigma$\cmmnt{;  yaitu, jika setiap subhimpunan terabaikan
dari $X$ terukur}.
     
\leader{211B}{Definisi} Misalkan $(X,\Sigma,\mu)$ suatu ruang ukur.
Maka $(X,\Sigma,\mu)$ adalah {\bf ruang peluang} jika $\mu X=1$.
Dalam hal ini $\mu$ disebut suatu {\bf peluang} atau {\bf ukuran
peluang}.
     
\leader{211C}{Definisi} Misalkan $(X,\Sigma,\mu)$ suatu ruang ukur.
Maka $\mu$, atau $(X,\Sigma,\mu)$, disebut {\bf hingga total} jika $\mu
X<\infty$.
     
\leader{211D}{Definisi} Misalkan $(X,\Sigma,\mu)$ suatu ruang ukur.
Maka $\mu$, atau $(X,\Sigma,\mu)$, disebut {\bf $\sigma$-hingga} jika
ada suatu barisan $\sequencen{E_n}$ yang terdiri atas himpunan-himpunan
terukur berukuran hingga sedemikian sehingga
$X=\bigcup_{n\in\Bbb N}E_n$.
     
\cmmnt{\medskip
     
\noindent{\bf Catatan} Perhatikan bahwa dalam hal ini kita dapat
menetapkan
     
\Centerline{$F_n=E_n\setminus\bigcup_{i<n}E_i$,\quad $G_n=\bigcup_{i\le
n}E_i$}
     
\noindent untuk setiap $n$, sehingga diperoleh suatu partisi
$\sequencen{F_n}$ dari $X$ (yaitu, suatu penutup saling lepas dari $X$)
menjadi himpunan-himpunan terukur berukuran hingga, dan suatu barisan
tak menurun $\sequencen{G_n}$ dari himpunan-himpunan berukuran hingga
yang menutupi $X$.
}%end of comment
     
\leader{211E}{Definisi} Misalkan $(X,\Sigma,\mu)$ suatu ruang ukur.
Maka $\mu$, atau $(X,\Sigma,\mu)$, disebut {\bf dapat dilokalkan secara
ketat} atau {\bf dapat didekomposisi} jika ada suatu partisi
$\langle X_i\rangle_{i\in I}$ dari $X$ menjadi himpunan-himpunan
terukur berukuran hingga sedemikian sehingga
     
\Centerline{$\Sigma
=\{E:E\subseteq X,\,E\cap X_i\in\Sigma\Forall i\in I\}$,}
     
\Centerline{$\mu E=\sum_{i\in I}\mu(E\cap X_i)$ untuk setiap $E\in\Sigma$.}
     
\noindent Saya akan menyebut keluarga demikian
$\langle X_i\rangle_{i\in I}$ suatu {\bf dekomposisi} dari $X$.
     
\cmmnt{\medskip
     
\noindent{\bf Catatan} Dalam konteks ini, kita dapat menafsirkan jumlah
$\sum_{i\in I}\mu(E\cap X_i)$ cukup sebagai
     
\Centerline{$\sup\{\sum_{i\in J}\mu(E\cap X_i):J$ adalah subhimpunan
hingga dari $I\}$,}
     
\noindent dengan mengambil
$\sum_{i\in\emptyset}\mu(E\cap X_i)=0$, karena kita hanya berurusan
dengan jumlah suku-suku nonnegatif (bdk.\ 112Bd).
}%end of comment
     
\leader{211F}{Definisi} Misalkan $(X,\Sigma,\mu)$ suatu ruang ukur.
Maka $\mu$, atau $(X,\Sigma,\mu)$, disebut {\bf semihingga} jika,
setiap kali $E\in\Sigma$ dan $\mu E=\infty$, ada suatu $F\subseteq E$
sedemikian sehingga $F\in\Sigma$ dan $0<\mu F<\infty$.
     
\leader{211G}{Definisi} Misalkan $(X,\Sigma,\mu)$ suatu ruang ukur.
Maka $\mu$, atau $(X,\Sigma,\mu)$, disebut {\bf dapat dilokalkan} jika
ia semihingga dan, untuk setiap $\Cal E\subseteq\Sigma$, ada suatu
$H\in\Sigma$ sedemikian sehingga (i) $E\setminus H$ terabaikan untuk
setiap $E\in\Cal E$ (ii) jika $G\in\Sigma$ dan $E\setminus G$
terabaikan untuk setiap $E\in\Cal E$, maka $H\setminus G$ terabaikan.
Akan berguna untuk menyebut himpunan demikian $H$ suatu {\bf supremum
esensial} dari $\Cal E$ dalam $\Sigma$.
     
\cmmnt{\medskip
     
\noindent{\bf Catatan} Definisi di sini agak canggung, karena
sebenarnya konsep ini berlaku bagi {\it aljabar} ukuran, bukan bagi
{\it ruang} ukur (lihat 211Ya-211Yb).   Namun, definisi yang sekarang
dapat dibuat berfungsi (lihat 213N, 241G, 243G di bawah) dan memungkinkan
kita melanjutkan tanpa pengantar formal terhadap konsep `aljabar
ukuran' sebelum tiba waktunya untuk mengerjakannya dengan semestinya
dalam Jilid 3.
}
     
\leader{211H}{Definisi} Misalkan $(X,\Sigma,\mu)$ suatu ruang ukur.
Maka $\mu$, atau $(X,\Sigma,\mu)$, disebut {\bf ditentukan secara
lokal} jika ia semihingga dan
     
\Centerline{$\Sigma=\{E:E\subseteq X,\,E\cap F\in\Sigma$ setiap kali
$F\in\Sigma$ dan $\mu F<\infty\}$\dvro{.}{;}}
     
\cmmnt{\noindent dengan kata lain, untuk setiap
$E\in\Cal PX\setminus\Sigma$ ada suatu $F\in\Sigma$ sedemikian sehingga
$\mu F<\infty$ dan $E\cap F\notin\Sigma$.}
     
\leader{211I}{Definisi} Misalkan $(X,\Sigma,\mu)$ suatu ruang ukur.
Suatu himpunan $E\in\Sigma$ adalah sebuah {\bf atom} bagi $\mu$ jika
$\mu E>0$ dan, setiap kali $F\in\Sigma$, $F\subseteq E$, salah satu
dari $F$, $E\setminus F$ terabaikan.
     
\leader{211J}{Definisi} Misalkan $(X,\Sigma,\mu)$ suatu ruang ukur.
Maka $\mu$, atau $(X,\Sigma,\mu)$, disebut {\bf tanpa atom} atau {\bf
terdifusi} jika tidak ada atom bagi $\mu$.
\cmmnt{(Perhatikan bahwa ini {\it tidak} sama dengan mengatakan bahwa
semua himpunan hingga terabaikan;  lihat 211R di bawah.   Sebagian
penulis menggunakan kata {\bf kontinu} dalam konteks ini.)}
     
\leader{211K}{Definisi} Misalkan $(X,\Sigma,\mu)$ suatu ruang ukur.
Maka $\mu$, atau $(X,\Sigma,\mu)$, disebut {\bf atomik murni} jika,
setiap kali $E\in\Sigma$ dan $E$ tidak terabaikan, ada suatu atom bagi
$\mu$ yang termuat dalam $E$.
     
\medskip
     
\noindent{\bf Catatan} \cmmnt{Ingat bahwa suatu ukuran $\mu$ pada
sebuah himpunan $X$ disebut {\bf bertumpu pada titik} jika $\mu$
mengukur setiap subhimpunan dari $X$ dan
$\mu E=\sum_{x\in E}\mu\{x\}$ untuk setiap $E\subseteq X$ (112Bd).}
Setiap ukuran yang bertumpu pada titik bersifat atomik murni\prooflet{,
karena $\{x\}$ mesti merupakan sebuah atom setiap kali
$\mu\{x\}>0$}, tetapi tidak setiap ukuran atomik murni bertumpu pada
titik\cmmnt{ (211R)}.
     
% {\smc Ramachandran 02} uses `discrete' for `purely atomic'
     
\leader{211L}{}\cmmnt{ Hubungan di antara konsep-konsep di atas dalam
suatu arti sangat langsung;  semua implikasi langsung ketika satu sifat
mengakibatkan sifat lain diberikan dalam teorema berikut.
     
\medskip
     
\noindent}{\bf Teorema} (a) Suatu ruang peluang bersifat
hingga total.
     
(b) Suatu ruang ukur hingga total bersifat $\sigma$-hingga.
     
(c) Suatu ruang ukur $\sigma$-hingga dapat dilokalkan secara ketat.
     
(d) Suatu ruang ukur yang dapat dilokalkan secara ketat dapat
dilokalkan dan ditentukan secara lokal.
     
(e) Suatu ruang ukur yang dapat dilokalkan bersifat semihingga.
     
(f) Suatu ruang ukur yang ditentukan secara lokal bersifat semihingga.
     
\proof{ (a), (b), (e), dan (f) langsung.
     
\medskip
     
{\bf (c)} Misalkan $(X,\Sigma,\mu)$ suatu ruang ukur $\sigma$-hingga;
misalkan $\sequencen{F_n}$ suatu barisan saling lepas yang terdiri atas
himpunan-himpunan terukur berukuran hingga dan menutupi $X$ (lihat
catatan dalam 211D).   Jika $E\in\Sigma$, tentu saja
$E\cap F_n\in\Sigma$ untuk setiap $n\in \Bbb N$, dan
     
\Centerline{$\mu E=\sum_{n=0}^{\infty}\mu(E\cap
F_n)=\sum_{n\in \Bbb N}\mu(E\cap F_n)$.}
     
\noindent Jika $E\subseteq X$ dan $E\cap F_n\in\Sigma$ untuk setiap
$n\in \Bbb N$, maka
     
\Centerline{$E=\bigcup_{n\in \Bbb N}E\cap F_n\in \Sigma$.}
     
\noindent Jadi $\sequencen{F_n}$ merupakan suatu dekomposisi dari $X$
dan $(X,\Sigma,\mu)$ dapat dilokalkan secara ketat.
     
\medskip
     
{\bf (d)} Misalkan $(X,\Sigma,\mu)$ suatu ruang ukur yang dapat
dilokalkan secara ketat;  misalkan $\langle X_i\rangle_{i\in I}$ suatu
dekomposisi dari $X$.
     
\medskip
     
\quad{\bf (i)} Misalkan $\Cal E$ suatu keluarga subhimpunan terukur dari
$X$.   Misalkan $\Cal F$ keluarga himpunan terukur $F\subseteq X$
sedemikian sehingga $\mu(F\cap E)=0$ untuk setiap $E\in\Cal E$.
Perhatikan bahwa $\emptyset\in\Cal F$ dan, jika $\sequencen{F_n}$ suatu
barisan sebarang dalam $\Cal F$, maka
$\bigcup_{n\in\Bbb N}F_n\in\Cal F$.   Untuk setiap $i\in I$, tetapkan
$\gamma_i=\sup\{\mu(F\cap X_i):F\in\Cal F\}$ dan pilih suatu barisan
$\sequencen{F_{in}}$ dalam $\Cal F$ sedemikian sehingga
$\lim_{n\to\infty}\mu(F_{in}\cap X_i)=\gamma_i$;  tetapkan
     
\Centerline{$F_i=\bigcup_{n\in\Bbb N}F_{in}\in\Cal F$.}
     
\noindent Tetapkan
     
\Centerline{$F=\bigcup_{i\in I}F_i\cap X_i\subseteq X$}
     
\noindent dan $H=X\setminus F$.
     
Kita melihat bahwa $F\cap X_i=F_i\cap X_i$ untuk setiap $i\in I$
(karena $\langle X_i\rangle_{i\in I}$ saling lepas), sehingga
$F\in\Sigma$ dan $H\in \Sigma$.   Untuk sebarang $E\in\Cal E$,
     
\Centerline{$\mu(E\setminus H)=\mu(E\cap F)
=\sum_{i\in I}\mu(E\cap F\cap X_i)
=\sum_{i\in I}\mu(E\cap F_i\cap X_i)=0$}
     
\noindent karena setiap $F_i$ termasuk dalam $\Cal F$.   Jadi
$F\in\Cal F$.   Jika $G\in\Sigma$ dan
$\mu(E\setminus G)=0$ untuk setiap $E\in\Cal E$, maka $X\setminus G$
dan $F'=F\cup(X\setminus G)$ termasuk dalam $\Cal F$.   Maka
$\mu(F'\cap X_i)\le\gamma_i$ untuk setiap $i\in I$.   Tetapi juga
$\mu(F\cap X_i)\ge\sup_{n\in\Bbb N}\mu(F_{in}\cap X_i)=\gamma_i$,
sehingga $\mu(F\cap X_i)=\mu(F'\cap X_i)$ untuk setiap $i$.
{\it Karena $\mu X_i$ hingga,} diperoleh
$\mu((F'\setminus F)\cap X_i)=0$ untuk setiap $i$.   Dengan menjumlahkan
atas $i$, $\mu(F'\setminus F)=0$, yaitu,
$\mu(H\setminus G)=0$.
     
Dengan demikian $H$ merupakan supremum esensial bagi $\Cal E$ dalam
$\Sigma$.   Karena $\Cal E$ sebarang, $(X,\Sigma,\mu)$ dapat dilokalkan.
     
\medskip
     
\quad{\bf (ii)}
Jika $E\in\Sigma$ dan $\mu E=\infty$, maka ada suatu $i\in I$
sedemikian sehingga
     
\Centerline{$0<\mu(E\cap X_i)\le\mu X_i<\infty$;}
     
\noindent sehingga $(X,\Sigma,\mu)$ bersifat semihingga.   Jika
$E\subseteq X$ dan $E\cap F\in\Sigma$ setiap kali $\mu F<\infty$, maka,
khususnya, $E\cap X_i\in\Sigma$ untuk setiap $i\in I$, sehingga
$E\in\Sigma$;  jadi $(X,\Sigma,\mu)$ ditentukan secara lokal.
}%end of proof of 211L
     
\leader{211M}{Contoh:  ukuran Lebesgue}\cmmnt{ Mari kita meninjau ukuran
Lebesgue dari sudut pandang konsep-konsep di atas.}   Tuliskan $\mu$
untuk ukuran Lebesgue pada $\BbbR^r$\cmmnt{ dan $\Sigma$ untuk
domainnya}.
     
\spheader 211Ma $\mu$ lengkap\prooflet{, karena ia dikonstruksi dengan
metode \Caratheodory;  jika $A\subseteq E$ dan $\mu E=0$, maka
$\mu^*A=\mu^*E=0$ (dengan menuliskan $\mu^*$ untuk ukuran luar
Lebesgue), sehingga, untuk sebarang $B\subseteq\BbbR^r$,
     
\Centerline{$\mu^*(B\cap A)+\mu^*(B\setminus A)\le 0+\mu^*B=\mu^*B$,}
     
\noindent dan $A$ mesti terukur}.
     
\spheader 211Mb $\mu$ bersifat $\sigma$-hingga\prooflet{, karena
$\BbbR^r=\bigcup_{n\in\Bbb N}[-\tbf{n},\tbf{n}]$, dengan
$\tbf{n}$ menyatakan vektor $(n,\ldots,n)$, dan
$\mu[-\tbf{n},\tbf{n}]=(2n)^r<\infty$ untuk setiap $n$.
Tentu saja $\mu$ bukan hingga total dan juga bukan ukuran peluang}.
     
\spheader 211Mc\cmmnt{ Karena} $\mu$
\cmmnt{bersifat $\sigma$-hingga, ia} dapat dilokalkan secara
ketat\cmmnt{ (211Lc)}, dapat dilokalkan\cmmnt{ (211Ld)}, ditentukan
secara lokal\cmmnt{ (211Ld)}, dan semihingga\cmmnt{ (211Le atau 211Lf)}.
     
\spheader 211Md $\mu$ tanpa atom.   \prooflet{\Prf\ Misalkan
$E\in\Sigma,\ \mu E>0$.  Tinjau fungsi
     
\Centerline{$a\mapsto f(a)
=\mu(E\cap[-\tbf{a},\tbf{a}]):\coint{0,\infty}\to\Bbb R$}
     
\noindent Kita mempunyai
     
\Centerline{$f(a)
\le f(b)\le f(a)+\mu[-\tbf{b},\tbf{b}]-\mu[-\tbf{a},\tbf{a}]
=f(a)+(2b)^r-(2a)^r$}
     
\noindent setiap kali $a\le b$ dalam $\coint{0,\infty}$, sehingga $f$
kontinu.   Sekarang $f(0)=0$ dan
$\lim_{n\to\infty}f(n)=\mu E>0$.   Berdasarkan Teorema Nilai Antara ada
suatu $a\in\coint{0,\infty}$ sedemikian sehingga
$0<f(a)<\mu E$.   Jadi kita mempunyai
     
\Centerline{$0<\mu(E\cap[-\tbf{a},\tbf{a}])<\mu E$.}
     
\noindent Karena $E$ sebarang, $\mu$ tanpa atom.\ \Qed}
     
\header{211Me}{\bf (e)} Sekarang merupakan pengamatan langsung bahwa
$\mu$ tidak mungkin atomik murni\cmmnt{, karena $\BbbR^r$ sendiri
merupakan suatu himpunan berukuran positif yang tidak memuat atom apa
pun}.
     
\leader{211N}{Ukuran pencacahan} Ambil $X$ sebagai sebarang himpunan tak
terhitung\cmmnt{ (misalnya, $X=\Bbb R$)}, dan $\mu$ sebagai ukuran
pencacahan pada $X$\cmmnt{ (112Bd)}.
     
\header{211Na}{\bf (a)} $\mu$
lengkap\ifwithproofs{, karena jika $A\subseteq E$ dan
$\mu E=0$ maka
     
\Centerline{$A=E=\emptyset\in\Sigma$.}
}\else{.}\fi
     
\header{211Nb}{\bf (b)} $\mu$ tidak $\sigma$-hingga\cmmnt{, karena jika
$\sequencen{E_n}$ suatu barisan sebarang dari himpunan-himpunan
berukuran hingga, maka setiap $E_n$ hingga, dan karena itu terhitung,
sedangkan $\bigcup_{n\in\Bbb N}E_n$ terhitung (1A1F), sehingga tidak
mungkin sama dengan $X$}.   \cmmnt{{\it Terlebih lagi,}} $\mu$ bukan
ukuran peluang dan juga bukan hingga total.
     
\spheader 211Nc $\mu$ dapat dilokalkan secara ketat.
\prooflet{\Prf\ Tetapkan $X_x=\{x\}$ untuk setiap $x\in X$.  Maka
$\langle X_x\rangle_{x\in X}$ merupakan suatu partisi dari $X$, dan
untuk sebarang $E\subseteq X$
     
\Centerline{$\mu(E\cap X_x)=1$ jika $x\in E$,\quad$0$ jika tidak.}
     
\noindent Berdasarkan definisi $\mu$,
     
\Centerline{$\mu E=\sum_{x\in X}\mu(E\cap X_x)$}
     
\noindent untuk setiap $E\subseteq X$, dan
$\langle X_x\rangle_{x\in X}$ merupakan suatu dekomposisi dari
$X$.\ \Qed
     
Akibatnya} $\mu$ dapat dilokalkan, ditentukan secara lokal, dan
semihingga.
     
\header{211Nd}{\bf (d)} $\mu$ atomik murni.   \prooflet{\Prf\ $\{x\}$
merupakan suatu atom untuk setiap $x\in X$, dan jika $\mu E>0$, maka
tentu $E$ memuat $\{x\}$ untuk suatu $x$.\ \QeD}
\cmmnt{Jelas bahwa} $\mu$ tidak tanpa atom.
     
\leader{211O}{Ruang yang tidak semihingga} Tetapkan $X=\{0\}$,
$\Sigma=\{\emptyset,X\}$, $\mu\emptyset=0$, dan $\mu X=\infty$.   Maka
$\mu$ tidak semihingga,\cmmnt{ karena $\mu X=\infty$ tetapi $X$ tidak
memiliki subhimpunan dengan ukuran hingga yang tak nol.   Akibatnya
$\mu$ tidak mungkin} dapat dilokalkan, ditentukan secara lokal,
$\sigma$-hingga, hingga total, ataupun merupakan ukuran peluang.
\cmmnt{Karena $\Sigma=\Cal PX$,} $\mu$ lengkap.
\cmmnt{$X$ merupakan suatu atom bagi $\mu$, sehingga} $\mu$ atomik
murni (bahkan, ia bertumpu pada titik).
     
\leader{211P}{Ruang yang tidak lengkap} Tuliskan $\Cal B$ untuk
aljabar-sigma dari subhimpunan-subhimpunan Borel dalam
$\Bbb R$\cmmnt{ (111G)}, dan $\mu$ untuk pembatasan ukuran Lebesgue
pada $\Cal B$\cmmnt{ (ingat bahwa berdasarkan 114G setiap subhimpunan
Borel dari $\Bbb R$ terukur Lebesgue)}.
Maka $(\Bbb R,\Cal B,\mu)$ tanpa atom, $\sigma$-hingga, dan tidak
lengkap.
     
\proof{{\bf (a)} Untuk melihat bahwa $\mu$ tidak lengkap, ingat bahwa
ada suatu permutasi kontinu dan naik ketat
$g:[0,1]\to[0,1]$ sedemikian sehingga $\mu g[C]>0$, dengan $C$
menyatakan himpunan Cantor, sehingga ada suatu himpunan
$A\subseteq g[C]$ yang tidak terukur Lebesgue (134Ib).   Sekarang
$g^{-1}[A]\subseteq C$ tidak mungkin merupakan himpunan Borel, sebab
$\chi A=\chi(g^{-1}[A])\smallcirc g^{-1}$ tidak terukur Lebesgue, dan
karena itu tidak terukur Borel, sedangkan komposisi dua fungsi terukur
Borel adalah terukur Borel (121Eg);  jadi $g^{-1}[A]$ merupakan
subhimpunan tak terukur dari himpunan terabaikan $C$.
     
\medskip
     
{\bf (b)} Argumen selebihnya dalam 211M berlaku bagi $\mu$ sama baiknya
seperti bagi ukuran Lebesgue yang sebenarnya, sehingga $\mu$
$\sigma$-hingga dan tanpa atom.
}%end of proof of 211P
     
\cmmnt{\medskip
     
\noindent{\bf *Catatan} Argumen yang diberikan dalam (a) dapat
menimbulkan kesan yang sangat keliru.   Himpunan $A$ yang disebut di
sana hanya dapat dikonstruksi dengan menggunakan suatu bentuk kuat
aksioma pilihan.   Perangkat demikian tidak diperlukan untuk hasil di
sini.
Ada banyak metode untuk mengonstruksi subhimpunan non-Borel dari
himpunan Cantor, dan semuanya menerangi persoalan dengan cara yang
berbeda.   Tanpa bentuk apa pun dari aksioma pilihan, ada kesulitan
dengan konsep `himpunan Borel', serta kesulitan lain dengan konsep
`ukuran Lebesgue', yang akan saya bahas dalam Bab 56;  tetapi pilihan
terhitung sudah cukup untuk menjamin keberadaan suatu subhimpunan
non-Borel dari $\Bbb R$.   Untuk perincian mengenai salah satu
pendekatan yang mungkin, lihat 423L dalam Jilid 4.
}
     
\leader{211Q}{Beberapa ruang peluang} Dua konstruksi ruang
peluang\cmmnt{ yang jelas}\cmmnt{, dengan membatasi diri pada
metode-metode yang diuraikan dalam Jilid 1,} adalah
     
(a) ukuran subruang yang diinduksi oleh ukuran Lebesgue pada
$[0,1]$\cmmnt{ (131B)};
     
(b) ukuran yang bertumpu pada titik, yang diinduksi pada suatu himpunan
$X$ oleh suatu fungsi $h:X\to[0,1]$ sedemikian sehingga
$\sum_{x\in X}h(x)=1$, dengan menuliskan
$\mu E=\sum_{x\in E}h(x)$ untuk setiap $E\subseteq X$;  misalnya, jika
$X$ suatu himpunan tunggal $\{x\}$ dan $h(x)=1$, atau jika
$X=\Bbb N$ dan $h(n)=2^{-n-1}$.
     
Dari kedua konstruksi ini, (a) memberikan ukuran peluang tanpa atom
dan (b) memberikan ukuran peluang atomik murni.
     
\leader{211R}{Ukuran terhitung-koterhitung}\cmmnt{ Konstruksi berikut
adalah salah satu konstruksi dasar yang patut diingat ketika meninjau
ruang ukur abstrak.
     
\medskip
     
}{\bf (a)} Misalkan $X$ sebarang himpunan.
Misalkan $\Sigma$ keluarga himpunan $E\subseteq X$ sedemikian sehingga
$E$ atau $X\setminus E$ terhitung.   Maka $\Sigma$ merupakan suatu
aljabar-sigma dari subhimpunan-subhimpunan $X$.
\prooflet{\Prf\ {(i)} $\emptyset$ terhitung, sehingga termasuk dalam
$\Sigma$.  {(ii)} Syarat agar $E$ termasuk dalam $\Sigma$ simetris
antara $E$ dan $X\setminus E$, sehingga $X\setminus E\in\Sigma$ untuk
setiap $E\in\Sigma$.   {(iii)} Misalkan $\sequencen{E_n}$ suatu barisan
sebarang dalam $\Sigma$, dan tetapkan
$E=\bigcup_{n\in\Bbb N}E_n$.   Jika setiap $E_n$ terhitung, maka $E$
terhitung, sehingga termasuk dalam $\Sigma$.   Jika tidak, ada suatu
$n$ sedemikian sehingga $X\setminus E_n$ terhitung; dalam hal ini
$X\setminus E\subseteq X\setminus E_n$ terhitung, sehingga sekali lagi
$E\in\Sigma$.\ \QeD}   $\Sigma$ disebut {\bf aljabar-sigma
terhitung-koterhitung} dari $X$.
     
\header{211Rb}{\bf (b)} Sekarang tinjau fungsi
$\mu:\Sigma\to\{0,1\}$ yang didefinisikan dengan menetapkan $\mu E=0$
jika $E$ terhitung, dan $\mu E=1$ jika $E\in\Sigma$ dan $E$ tidak
terhitung.   Maka $\mu$ adalah suatu ukuran.
\prooflet{\Prf\ {(i)} $\emptyset$ terhitung, sehingga
$\mu\emptyset=0$.   {(ii)} Misalkan $\sequencen{E_n}$ suatu barisan
saling lepas dalam $\Sigma$, dan $E$ gabungannya.   {(${\alpha}$)}
Jika setiap $E_m$ terhitung, maka $E$ juga terhitung, sehingga
     
\Centerline{$\mu E=0=\sum_{n=0}^{\infty}\mu E_n$.}
     
\noindent{(${\beta}$)} Jika suatu $E_m$ tak terhitung, maka
$E\supseteq E_m$ juga tak terhitung, dan $\mu E=\mu E_m=1$.   Tetapi
dalam hal ini, karena $E_m\in\Sigma$, $X\setminus E_m$ terhitung,
sehingga $E_n$, sebagai suatu subhimpunan dari $X\setminus E_m$,
terhitung untuk setiap $n\ne m$; jadi $\mu E_n=0$ untuk setiap
$n\ne m$, dan
     
\Centerline{$\mu E=1=\sum_{n=0}^{\infty}\mu E_n$.}
     
\noindent Karena $\sequencen{E_n}$ sebarang, $\mu$ merupakan suatu
ukuran.\ \Qed}
Ini adalah {\bf ukuran terhitung-koterhitung} pada $X$.
     
\header{211Rc}{\bf (c)} Jika $X$ sebarang himpunan tak terhitung dan
$\mu$ ukuran terhitung-koterhitung pada $X$, maka $\mu$ merupakan
ukuran peluang lengkap dan atomik murni, tetapi tidak bertumpu pada
titik.   \prooflet{\Prf\ {(i)} Jika $A\subseteq E$ dan $\mu E=0$, maka
$E$ terhitung, sehingga $A$ juga terhitung dan termasuk dalam
$\Sigma$.   Jadi $\mu$ lengkap.   {(ii)} Karena $X$ tak terhitung,
$\mu X=1$ dan $\mu$ merupakan ukuran peluang.   {(iii)} Jika
$\mu E>0$, maka $\mu F=\mu E=1$ setiap kali $F$ merupakan subhimpunan
terukur yang tidak terabaikan dari $E$, sehingga $E$ sendiri merupakan
suatu atom; jadi $\mu$ atomik murni.   (iv)
$\mu X=1>0=\sum_{x\in X}\mu\{x\}$, sehingga $\mu$ tidak bertumpu pada
titik.\ \Qed}
     
\exercises{\leader{211X}{Latihan dasar $\pmb{>}$(a)}
%\sqheader 211Xa 
Misalkan $(X,\Sigma,\mu)$ suatu ruang ukur.   Tunjukkan bahwa $\mu$
$\sigma$-hingga jika dan hanya jika ada suatu ukuran hingga total
$\nu$ pada $X$ yang memiliki himpunan-himpunan terukur dan
himpunan-himpunan terabaikan yang sama dengan $\mu$.
%211D
     
\sqheader 211Xb Misalkan $g:\Bbb R\to\Bbb R$ suatu fungsi tak menurun
dan $\mu_g$ ukuran Lebesgue-Stieltjes yang bersesuaian (114Xa).
Tunjukkan bahwa $\mu_g$ lengkap dan $\sigma$-hingga.   Tunjukkan bahwa
     
\quad (i) $\mu_g$ hingga total jika dan hanya jika $g$ terbatas;
     
\quad (ii) $\mu_g$ merupakan ukuran peluang jika dan hanya jika
$\lim_{x\to\infty}g(x) -\lim_{x\to-\infty}g(x)=1$;
     
\quad (iii) $\mu_g$ tanpa atom jika dan hanya jika $g$ kontinu;
     
\quad (iv) jika $E$ sebarang atom bagi $\mu_g$, ada suatu titik
$x\in E$ sedemikian sehingga $\mu_gE=\mu_g\{x\}$;
     
\quad (v) $\mu_g$ atomik murni jika dan hanya jika ia bertumpu pada
titik.
%211M
     
\sqheader 211Xc Misalkan $\mu$ ukuran pencacahan pada suatu himpunan
$X$.   Tunjukkan bahwa $\mu$ selalu lengkap, dapat dilokalkan secara
ketat, dan atomik murni, serta bahwa ia $\sigma$-hingga jika dan hanya
jika $X$ terhitung, hingga total jika dan hanya jika $X$ hingga,
merupakan ukuran peluang jika dan hanya jika $X$ suatu himpunan
tunggal, dan tanpa atom jika dan hanya jika $X$ kosong.
%211N
     
\spheader 211Xd Tunjukkan bahwa suatu ukuran yang bertumpu pada titik
selalu lengkap, dan dapat dilokalkan secara ketat jika dan hanya jika
ia semihingga.
%211O
     
\spheader 211Xe Misalkan $X$ suatu himpunan.   Tunjukkan bahwa untuk
sebarang ideal-$\sigma$ $\Cal I$ dari subhimpunan-subhimpunan $X$
(definisi:  112Db), himpunan
     
\Centerline{$\Sigma=\Cal I\cup\{X\setminus A:A\in\Cal I\}$}
     
\noindent merupakan suatu aljabar-sigma dari subhimpunan-subhimpunan
$X$, dan bahwa ada suatu ukuran $\mu:\Sigma\to\{0,1\}$ yang diberikan
dengan menetapkan
     
\Centerline{$\mu E=0$ jika $E\in\Cal I$,
\quad $\mu E=1$ jika $E\in\Sigma\setminus\Cal I$.}
     
\noindent Tunjukkan bahwa $\Cal I$ tepat merupakan ideal nol dari
$\mu$, bahwa $\mu$ lengkap, hingga total, dan atomik murni, serta
merupakan ukuran peluang jika dan hanya jika $X\notin\Cal I$.
%211R
     
\spheader 211Xf Tunjukkan bahwa suatu ukuran yang bertumpu pada titik
dapat dilokalkan secara ketat jika dan hanya jika ia semihingga.
     
\spheader 211Xg Misalkan $(X,\Sigma,\mu)$ suatu ruang ukur sedemikian
sehingga $\mu X>0$.   Tunjukkan bahwa himpunan semua subhimpunan
koterabaikan dari $X$ merupakan suatu filter pada $X$.

\leader{211Y}{Latihan lanjutan (a)}
%\spheader 211Ya
Misalkan $(X,\Sigma,\mu)$ suatu ruang ukur, dan untuk
$E$, $F\in\Sigma$ tuliskan $E\sim F$ jika
$\mu(E\symmdiff F)=0$.   Tunjukkan bahwa $\sim$ merupakan suatu relasi
ekuivalensi pada $\Sigma$.   Misalkan $\frak A$ himpunan kelas-kelas
ekuivalensi dalam $\Sigma$ untuk $\sim$;  untuk $E\in\Sigma$, tuliskan
$E^{\ssbullet}\in\frak A$ untuk kelas ekuivalensinya.   Tunjukkan bahwa
ada suatu relasi urutan parsial $\Bsubseteq$ pada $\frak A$ yang
didefinisikan dengan mengatakan bahwa, untuk $E$, $F\in\Sigma$,
     
\Centerline{$E^{\ssbullet}\Bsubseteq F^{\ssbullet}
\,\iff\,\mu(E\setminus F)=0$.}
     
\noindent Tunjukkan bahwa $\mu$ dapat dilokalkan jika dan hanya jika,
untuk setiap $A\subseteq\frak A$, ada suatu $h\in\frak A$ sedemikian
sehingga (i) $a\Bsubseteq h$ untuk setiap $a\in A$ (ii) setiap kali
$g\in\frak A$ sedemikian sehingga $a\Bsubseteq g$ untuk setiap
$a\in A$, maka $h\Bsubseteq g$.
%211G
     
\spheader 211Yb Misalkan $(X,\Sigma,\mu)$ suatu ruang ukur, dan
konstruksikan $\frak A$ seperti dalam 211Ya.   Tunjukkan bahwa ada
operasi-operasi $\Bcup$, $\Bcap$, $\Bsetminus$ pada $\frak A$ yang
didefinisikan dengan mengatakan bahwa
     
\Centerline{$E^{\ssbullet}\Bcap F^{\ssbullet}=(E\cap F)^{\ssbullet}$,
\quad$E^{\ssbullet}\Bcup F^{\ssbullet}=(E\cup F)^{\ssbullet}$,
\quad$E^{\ssbullet}\Bsetminus F^{\ssbullet}=(E\setminus F)^{\ssbullet}$}
     
\noindent untuk semua $E$, $F\in\Sigma$.
Tunjukkan bahwa jika $A\subseteq\frak A$ sebarang himpunan
{\it terhitung}, maka ada suatu $h\in\frak A$ sedemikian sehingga
(i) $a\Bsubseteq h$ untuk setiap $a\in A$ (ii) setiap kali
$g\in\frak A$ sedemikian sehingga $a\Bsubseteq g$ untuk setiap
$a\in A$, maka $h\Bsubseteq g$.   Tunjukkan bahwa ada suatu
fungsional $\bar\mu:\frak A\to[0,\infty]$ yang didefinisikan dengan
mengatakan bahwa $\bar\mu(E^{\ssbullet})=\mu E$ untuk setiap
$E\in\Sigma$.
($(\frak A,\bar\mu)$ disebut {\bf aljabar ukuran} dari
$(X,\Sigma,\mu)$.)
%211Ya 211G
     
\spheader 211Yc Misalkan $(X,\Sigma,\mu)$ suatu ruang ukur semihingga.
Tunjukkan bahwa ruang ini tanpa atom jika dan hanya jika, setiap kali
$\epsilon>0$, $E\in\Sigma$, dan $\mu E<\infty$, ada suatu partisi
hingga dari $E$ menjadi himpunan-himpunan terukur yang masing-masing
berukuran paling besar $\epsilon$.
%211J
     
\spheader 211Yd Misalkan $(X,\Sigma,\mu)$ suatu ruang ukur yang dapat
dilokalkan secara ketat.   Tunjukkan bahwa ruang ini tanpa atom jika dan
hanya jika, untuk setiap $\epsilon>0$, ada suatu dekomposisi dari $X$
yang terdiri atas himpunan-himpunan berukuran paling besar $\epsilon$.
%211Yc 211J
     
\spheader 211Ye Misalkan $\Sigma$ aljabar-sigma
terhitung-koterhitung dari $\Bbb R$.   Tunjukkan bahwa
$\coint{0,\infty}\notin\Sigma$.   Misalkan $\mu$ pembatasan ukuran
pencacahan pada $\Sigma$.   Tunjukkan bahwa $(\Bbb R,\Sigma,\mu)$
lengkap, semihingga, dan atomik murni, tetapi tidak dapat dilokalkan dan
tidak ditentukan secara lokal.
%211R
}%end of exercises
     
\endnotes{
\Notesheader{211} Daftar definisi dalam 211A-211K mungkin terasa sudah
cukup panjang bagi Anda, sekalipun saya telah menghilangkan banyak
gagasan yang kadang-kadang berguna.   Konsep-konsep di sini sangat
beragam dalam hal kepentingannya, dan kepentingan masing-masing juga
sangat bergantung pada konteks.   Pandangan saya sendiri adalah bahwa,
ketika mempelajari sebarang ruang ukur, kita mutlak perlu mengetahui
klasifikasinya menurut sebelas ciri pembeda yang dicantumkan di sini,
dan mampu menguraikan setiap atom yang ada.   Untungnya, untuk
kebanyakan ruang ukur `biasa', klasifikasi itu cukup cepat, karena
jika (misalnya) ruang tersebut $\sigma$-hingga dan Anda mengetahui
ukuran seluruh ruang, satu-satunya pertanyaan yang tersisa menyangkut
kelengkapan dan atom.   Perbedaan antara ruang-ruang yang dapat atau
tidak dapat dilokalkan secara ketat, semihingga, dapat dilokalkan, dan
ditentukan secara lokal hanya relevan bagi ruang-ruang yang tidak
$\sigma$-hingga, dan tidak muncul dalam penerapan elementer.
     
Menurut saya, wajar pula untuk mengatakan bahwa gagasan ruang ukur
`lengkap' dan `ditentukan secara lokal' bersifat teknis;  maksud
saya, gagasan-gagasan itu tidak bersesuaian dengan ciri penting dalam
struktur esensial suatu ruang, walaupun ada beberapa masalah menarik
mengenai ukuran yang tidak lengkap.   Salah satu perwujudannya adalah
adanya konstruksi kanonik untuk membuat ruang menjadi lengkap, atau
lengkap sekaligus ditentukan secara lokal (212C, 213D-213E).   Selain
itu, ruang ukur yang tidak semihingga sesungguhnya tidak termasuk dalam
teori ukuran, melainkan dalam kajian yang lebih umum mengenai
aljabar-sigma dan ideal-$\sigma$.   Klasifikasi yang paling penting,
ditinjau dari perilaku suatu ruang ukur, menurut saya adalah
`$\sigma$-hingga', `dapat dilokalkan', dan `dapat
dilokalkan secara ketat';  inilah ciri-ciri kritis yang tidak dapat
dipaksakan melalui konstruksi elementer.
     
Jika Anda mengetahui sesuatu tentang subhimpunan Borel dari garis real,
argumen dalam bagian (a) dari bukti 211P pasti tampak sangat canggung.
Namun, bukti yang `lebih baik' bergantung pada gagasan yang baru akan
kita perlukan dalam Jilid 4, sedangkan bukti di sini didasarkan pada
suatu konstruksi yang harus kita pahami untuk alasan-alasan lain.
}%end of notes
     
\discrpage
    
     




